\documentclass[11pt,a4paper]{article}

% ── Packages ──────────────────────────────────────────────────────────────────
\usepackage[margin=1in]{geometry}
\usepackage{amsmath,amssymb}
\usepackage{booktabs}
\usepackage{graphicx}
\usepackage{caption}
\usepackage{subcaption}
\usepackage{hyperref}
\usepackage{xcolor}
\usepackage{array}
\usepackage{multirow}
\usepackage{enumitem}
\usepackage{fancyhdr}
\usepackage{titlesec}
\usepackage{parskip}
\setlength{\emergencystretch}{3em}
\usepackage[T1]{fontenc}
\usepackage{lmodern}
\usepackage{listings}
\usepackage{inconsolata}

% ── Listings style (Python) ───────────────────────────────────────────────────
\lstset{
  language=Python,
  basicstyle=\footnotesize\ttfamily,
  keywordstyle=\color{arbblue}\bfseries,
  stringstyle=\color{manipred},
  commentstyle=\color{gray}\itshape,
  numberstyle=\tiny\color{gray},
  numbers=left,
  numbersep=6pt,
  stepnumber=5,
  frame=single,
  framerule=0.4pt,
  breaklines=true,
  breakatwhitespace=false,
  showstringspaces=false,
  tabsize=4,
  captionpos=t,
  aboveskip=8pt,
  belowskip=4pt,
  xleftmargin=14pt,
  xrightmargin=0pt,
  morekeywords={True,False,None,self},
}

% ── Header / Footer ───────────────────────────────────────────────────────────
\pagestyle{fancy}
\fancyhf{}
\rhead{\small ORIE5570 --- Part 3}
\lhead{\small Ishan Patwardhan (ip259)}
\cfoot{\thepage}
\renewcommand{\headrulewidth}{0.4pt}

% ── Section formatting ────────────────────────────────────────────────────────
\titleformat{\section}{\large\bfseries}{{\thesection.}}{0.5em}{}[\titlerule]
\titleformat{\subsection}{\normalsize\bfseries}{{\thesubsection}}{0.5em}{}

% ── Colours ──────────────────────────────────────────────────────────────────
\definecolor{arbblue}{HTML}{2166AC}
\definecolor{manipred}{HTML}{D6604D}
\definecolor{herdorange}{HTML}{E08214}

% ── Hyperref ─────────────────────────────────────────────────────────────────
\hypersetup{colorlinks=true, linkcolor=black, citecolor=black, urlcolor=arbblue}

% ─────────────────────────────────────────────────────────────────────────────
\begin{document}

% ── Title block ──────────────────────────────────────────────────────────────
\begin{center}
  {\LARGE\bfseries Part 3: REINFORCE with Baseline}\\[4pt]
  {\large Multi-Trader-Type Crypto Market Analysis}\\[8pt]
  {\large BTCB-USD (Bitcoin BEP2) $\mid$ Policy Gradient vs.\ DQN}\\[12pt]
  \begin{tabular}{rl}
    \textbf{Name:}   & Ishan Patwardhan \\
    \textbf{NetID:}  & ip259 \\
    \textbf{Course:} & ORIE5570 --- Reinforcement Learning with Operations Research Applications \\
  \end{tabular}
\end{center}
\vspace{-4pt}\rule{\linewidth}{0.8pt}\vspace{6pt}

% ─────────────────────────────────────────────────────────────────────────────
\section{Methodology}

\subsection{REINFORCE with Baseline}

REINFORCE is an on-policy Monte Carlo policy gradient algorithm.  Unlike DQN, which learns a deterministic Q-function off-policy via experience replay, REINFORCE directly parameterises and optimises the policy $\pi(a\mid s;\theta)$ using the policy gradient theorem:
\[
  \nabla_\theta J(\theta)
  = \mathbb{E}\!\left[\,\sum_t \nabla_\theta \log \pi(a_t\mid s_t;\theta)
    \cdot \bigl(G_t - b(s_t)\bigr)\right],
\]
where $G_t = \sum_{k\ge t}\gamma^{k-t}r_k$ is the discounted return and $b(s_t)$ is the \textbf{baseline} --- a learned state-value function $V(s;\phi)$ that reduces gradient variance without introducing bias.

\medskip
\textbf{Architecture.}  Both networks are \textbf{two-layer MLPs} with ReLU activations and 128 hidden units per layer:

\begin{table}[h!]
\centering
\small
\begin{tabular}{lccccl}
\toprule
\textbf{Network} & \textbf{Input} & \textbf{Hidden 1} & \textbf{Hidden 2} & \textbf{Output} & \textbf{Final act.}\\
\midrule
Policy $\pi(a\mid s;\theta)$ & 13 & 128, ReLU & 128, ReLU & 3 logits & Softmax\\
Baseline $V(s;\phi)$        & 13 & 128, ReLU & 128, ReLU & 1 scalar  & Linear\\
\bottomrule
\end{tabular}
\end{table}

\textbf{Training protocol.}  Each episode spans 168 steps (one week of hourly data).  The agent collects a full trajectory using its stochastic policy, computes discounted returns $G_t$, normalises the advantage $A_t = G_t - V(s_t)$, and updates both networks with Adam ($\text{lr}=3\!\times\!10^{-4}$, $\gamma=0.99$, gradient clipping at~1.0).  Each trader type is trained for \textbf{1{,}200 episodes} on the training set (March~2024 -- October~2025).

\subsection{Three Trader Types}

Three economically distinct agents are modelled by shaping the reward on top of the base log-return signal $r_t = \mathrm{clip}\!\left(\log(V_t/V_{t-1}),\,-1,\,+1\right)$:

\begin{table}[h!]
\centering
\small
\renewcommand{\arraystretch}{1.3}
\begin{tabular}{>{\raggedright}p{2.8cm} >{\raggedright}p{5.2cm} >{\raggedright\arraybackslash}p{5.5cm}}
\toprule
\textbf{Trader Type} & \textbf{Reward Modification} & \textbf{Economic Rationale}\\
\midrule
\textcolor{arbblue}{\textbf{Rational Arbitrageur}} &
  $r - 0.4\cdot\sigma_{24}\cdot f_{\text{crypto}}$ &
  Penalises variance of crypto exposure; prefers risk-adjusted returns.  Models stat-arb agents who close positions quickly to limit overnight vol.\\
\textcolor{manipred}{\textbf{Manipulator}} (front-runner) &
  $r + 0.15\cdot\bar{V}\cdot\mathrm{sign}(r_{4h})\cdot d$ &
  Bonus for trading in the direction of large-volume moves.  $\bar{V}$ = standardised log-volume; $d\in\{+1,0,-1\}$.  Models front-runners who detect institutional order flow via volume spikes.\\
\textcolor{herdorange}{\textbf{Herd-following Retail}} &
  $r + 0.25\cdot\mathrm{sign}(r_{24h})\cdot d$ &
  Bonus for aligning with the 24-hour trend.  Models FOMO-driven retail who chase recent price direction regardless of risk.\\
\bottomrule
\end{tabular}
\end{table}

\subsection{Shapley Value Computation}

Shapley values from cooperative game theory are the unique attribution satisfying efficiency, symmetry, dummy, and additivity axioms. For each trained policy network, SHAP's \texttt{KernelExplainer} is applied to the Buy-action probability $\pi(\text{Buy}\mid s;\theta)$, treating the network as a black box. Background: 30 test states; evaluation: 200 test states, 80 coalition samples each. Mean absolute SHAP $|\phi_i|$ --- the average magnitude of feature~$i$'s contribution to buy-action probability --- is reported for all 12 features.

\subsection{Market Simulation}

\sloppy
To measure systemic effects, a simulation is run over \textbf{500 consecutive test-set hours} under 7 trader compositions $(w_{\text{arb}},w_{\text{man}},w_{\text{herd}})$. At each step the aggregate signal $S_t = \sum_k w_k\,d_k$, $d_k\in\{+1,0,-1\}$. Three metrics are recorded:
\fussy

\begin{itemize}[nosep]
  \item \textbf{Crisis frequency:} fraction of steps where regime\,=\,\texttt{decline} and $S_t < -0.3$ simultaneously.
  \item \textbf{Stablecoin stability:} USDT-USD hourly peg deviation data (Yahoo Finance, 3{,}394 hours, mean dev.\,=\,0.0576\%). Average USDT deviation from \$1.00 during $|S_t|>0.3$ periods; metric = $1/\text{(mean deviation)}$, higher\,=\,more stable.
  \item \textbf{Market efficiency:} $1 - |\rho_1|$, where $\rho_1$ is the lag-1 autocorrelation of simulated hourly returns.
\end{itemize}

% ─────────────────────────────────────────────────────────────────────────────
\section{Results}

\subsection{Backtest Performance Comparison}

All agents are evaluated on the out-of-sample test period (October~2025 -- February~2026), during which BTCB-USD declined 46.2\%.

\begin{table}[h!]
\centering
\caption{Backtest Performance --- All Agents (Test Period Oct~2025 -- Feb~2026)}
\small
\begin{tabular}{p{4.2cm} p{2.8cm} r r r r}
\toprule
\textbf{Agent} & \textbf{Method} & \textbf{Return} & \textbf{Sharpe} & \textbf{Max DD} & \textbf{Final PV}\\
\midrule
Buy \& Hold         & Passive            & $-46.21\%$ & $-3.201$ & $-49.99\%$ & \$5{,}379\\
DQN (Part~2)        & Off-policy, det.   & $-49.44\%$ & $-5.301$ & $-53.18\%$ & \$5{,}056\\
REINF.\ \textcolor{arbblue}{Arbitrageur}
                    & On-policy, stoch.  & $-53.07\%$ & $-5.424$ & $-56.23\%$ & \$4{,}693\\
REINF.\ \textcolor{herdorange}{Herder}
                    & On-policy, stoch.  & $-17.24\%$ & $\mathbf{-1.550}$ & $-26.32\%$ & \$8{,}276\\
\textbf{REINF.\ \textcolor{manipred}{Manipulator}}
                    & On-policy, stoch.  & $\mathbf{-6.01\%}$ & $-3.936$ & $\mathbf{-6.23\%}$ & $\mathbf{\$9{,}399}$\\
\bottomrule
\end{tabular}
\end{table}

\sloppy
\textbf{Key findings.} The manipulator REINFORCE agent preserved \textbf{94\% of capital} during a 46\% market crash. Its front-running reward developed a highly selective buy trigger: only entering when large volume \emph{and} positive 4-hour momentum coincide --- a signal that rarely fires in a bear market. The arbitrageur performed worst ($-53\%$): its volatility penalty caused it to hold large positions during rallies then sell aggressively as volatility spiked.

\textbf{REINFORCE vs.\ DQN.} REINFORCE's stochastic policy proved more adaptive than DQN's deterministic convergence (Q-values near-uniform: Hold\,=\,0.0043, Buy\,=\,0.0041, Sell\,=\,0.0033). REINFORCE maintains exploration throughout deployment; DQN's $\arg\max$ collapses once Q-values converge. However, reward design dominates architecture: the arbitrageur REINFORCE \emph{underperforms} DQN, confirming that reward shaping matters more than the choice between policy-gradient and value-based methods.

Figure~\ref{fig:training} shows the 100-episode rolling average reward for all three REINFORCE agents during training.

\begin{figure}[h!]
  \centering
  \includegraphics[width=\linewidth]{part3_training.png}
  \caption{REINFORCE training curves (100-episode rolling average).  The manipulator agent converges to positive average reward; the herder shows high variance consistent with its momentum-following reward shaping; the arbitrageur plateaus at a negative reward due to its strict volatility penalty.}
  \label{fig:training}
\end{figure}

\subsection{Shapley Value Analysis --- Feature Importance by Trader Type}

\begin{table}[h!]
\centering
\caption{Mean Absolute SHAP Values for Buy-Action Probability (all 12 features)}
\small
\begin{tabular}{clclcl}
\toprule
\multicolumn{2}{c}{\textcolor{arbblue}{\textbf{Arbitrageur}}} &
\multicolumn{2}{c}{\textcolor{manipred}{\textbf{Manipulator}}} &
\multicolumn{2}{c}{\textcolor{herdorange}{\textbf{Herder}}}\\
\cmidrule(lr){1-2}\cmidrule(lr){3-4}\cmidrule(lr){5-6}
Feature & $|\phi|$ & Feature & $|\phi|$ & Feature & $|\phi|$\\
\midrule
\textbf{rsi\_14}        & 0.01595 & \textbf{price\_vs\_ma50} & 0.02006 & \textbf{ret\_168h}      & 0.01608\\
\textbf{ret\_24h}       & 0.01310 & \textbf{vol\_24h}        & 0.01612 & \textbf{price\_vs\_ma20}& 0.01188\\
\textbf{hour\_cos}      & 0.01299 & \textbf{price\_vs\_ma20} & 0.01437 & \textbf{price\_vs\_ma50}& 0.01171\\
ret\_168h               & 0.01243 & log\_vol                 & 0.01224 & rsi\_14                 & 0.01122\\
price\_vs\_ma20         & 0.01198 & hl\_range                & 0.01196 & vol\_24h                & 0.01095\\
\bottomrule
\end{tabular}
\end{table}

Figure~\ref{fig:shap} shows the full SHAP feature importance profiles for all 12 features across trader types.

\begin{figure}[h!]
  \centering
  \includegraphics[width=\linewidth]{part3_shap.png}
  \caption{Mean $|\text{SHAP}|$ for Buy-action probability across all 12 market features, for each trader type.  Feature importance profiles are distinct and economically interpretable.}
  \label{fig:shap}
\end{figure}

\textbf{Top-3 feature explanations by trader type:}

\begin{description}[leftmargin=1.8cm, labelwidth=1.6cm, style=nextline]
  \item[\textcolor{arbblue}{Arb:} \texttt{rsi\_14}, \texttt{ret\_24h}, \texttt{hour\_cos}]
    RSI ($|\phi|=0.016$) is the arbitrageur's primary signal: extreme overbought/oversold readings identify mean-reversion opportunities.  The 24-hour return acts as a regime filter --- the agent avoids buying into sustained downtrends where mean-reversion takes too long to materialise.  Hour-of-day cosine captures circadian liquidity; the arbitrageur avoids thin Asian overnight markets because its volatility penalty makes illiquid execution costly.

  \item[\textcolor{manipred}{Manip:} \texttt{price\_vs\_ma50}, \texttt{vol\_24h}, \texttt{price\_vs\_ma20}]
    Long-term MA deviation ($|\phi|=0.020$) is the front-runner's highest-weight feature: large investors typically buy/sell systematically as price recovers toward trend, giving the manipulator a predictable order-flow target.  Realised volatility ($|\phi|=0.016$) is the trigger: high vol signals active institutional participation and therefore large orders to front-run.  Critically, the top-5 features are \emph{all} microstructure signals rather than macroeconomic trend signals --- consistent with how real front-runners operate.

  \item[\textcolor{herdorange}{Herd:} \texttt{ret\_168h}, \texttt{price\_vs\_ma20}, \texttt{price\_vs\_ma50}]
    The 7-day return ($|\phi|=0.016$) is dominant: retail herders follow sustained multi-day trends, not hourly signals.  Both MA deviation features confirm the trend direction before the herder commits.  Notably absent from the herder top-5: RSI and volume --- herders are blind to contrarian signals and institutional order flow, making them systematically exploitable by both arbitrageurs and manipulators.
\end{description}

\subsection{Regulatory Implications and Market Simulation}

\begin{table}[h!]
\centering
\caption{Market Metrics Under Different Trader Compositions (500-step simulation)}
\small
\begin{tabular}{lrrr}
\toprule
\textbf{Composition} & \textbf{Crisis Freq.} & \textbf{Stablecoin Stability} & \textbf{Market Efficiency}\\
\midrule
All Arbitrageurs     & 0.008  & 18.14 & 0.871\\
All Manipulators     & \textbf{0.000}  & \textbf{22.57} & 0.865\\
All Herders          & 0.054  & 21.00 & 0.870\\
Equal Mix (33/33/34) & 0.034  & 21.82 & 0.869\\
Arb-Dominated (60/20/20) & 0.008  & 18.14 & 0.870\\
Manip-Dominated (20/60/20) & 0.004  & 21.25 & 0.867\\
\textbf{Herd-Dominated (20/20/60)} & \textbf{0.054}  & 21.00 & \textbf{0.869}\\
\bottomrule
\end{tabular}
\small\textit{Stablecoin Stability = $1/(\text{mean USDT peg deviation during }|S_t|>0.3)$; higher = more stable.}
\end{table}

Figure~\ref{fig:market} shows all three metrics across the seven compositions.

\begin{figure}[h!]
  \centering
  \includegraphics[width=\linewidth]{part3_market_metrics.png}
  \caption{Crisis frequency (left), stablecoin stability (centre), and market efficiency (right) under seven trader-type compositions.  Dashed horizontal line marks the equal-mix baseline.  Arbitrageur-dominated markets stress the USDT peg most (lowest stability); herd-dominated markets generate the most crisis events.}
  \label{fig:market}
\end{figure}

\textbf{Four regulatory insights:}

\begin{enumerate}[leftmargin=*,nosep]
  \item \textbf{Arbitrageurs maximally stress stablecoin peg.}  Arbitrageur-dominated markets produce the lowest stablecoin stability (18.14) because arbitrageurs generate concentrated sell signals exactly when stablecoin redemption pressure peaks.  Regulators should treat a surge in arbitrageur-style algorithmic selling as an early warning of peg stress.

  \item \textbf{Manipulators paradoxically protect stablecoin stability.}  Manipulator-dominated markets show the highest stability (22.57) and zero crisis events.  This is a false negative: front-runners smooth abrupt cascades by acting before price breaks, masking wealth extraction behind benign aggregate metrics.  Regulators must use order-book surveillance (clusters of pre-large-order positions), not price-level monitoring, to detect manipulation.

  \item \textbf{Herders generate the highest crisis frequency.}  Herd-dominated markets produce 6.75$\times$ more crisis events (0.054) than arbitrageur markets (0.008), amplifying sell signals during confirmed downtrends into positive feedback loops.  Stablecoin stability is moderate (21.00) because herder crises are gradual, reducing concentrated redemption pressure relative to sharp arbitrageur events.

  \item \textbf{Equal-mix composition is the regulatory sweet spot.}  The equal mix (33/33/34) achieves a balanced profile: low crisis frequency (0.034), good stablecoin stability (21.82), and high efficiency (0.869).  Policy should attract arbitrageur activity in calm periods (improving efficiency) while dampening momentum orders during declining markets through friction rules that apply asymmetrically to trend-following but not mean-reversion orders.
\end{enumerate}

% ─────────────────────────────────────────────────────────────────────────────
\section{Conclusions}

This study extended the Part~2 BTCB-USD RL framework in three directions.  First, \textbf{REINFORCE with baseline} was implemented as a stochastic on-policy policy gradient agent.  Its distributional policy proved more adaptive than DQN's deterministic cash-hold convergence, with the manipulator-type REINFORCE preserving 94\% of capital ($-6\%$ vs.\ $-46\%$ buy-and-hold) through a highly selective entry condition grounded in volume and momentum signals.  Second, \textbf{Shapley value attribution} revealed that the three trader types use fundamentally different information channels: arbitrageurs rely on mean-reversion indicators (RSI, 24h return), manipulators exploit microstructure signals (long-term MA deviation, realised volatility, volume), and herders weight long-horizon trend signals (7-day return, MA crossovers).  These profiles are directly actionable for surveillance system design.  Third, \textbf{multi-agent market simulation} grounded in real USDT-USD peg data demonstrated that arbitrageur concentration maximally stresses stablecoin peg stability, herder concentration maximally inflates crisis frequency, and a balanced trader composition achieves the best regulatory outcome across all three metrics.

\vspace{6pt}
\noindent\small\textit{All code: \texttt{rl\_crypto\_part3.ipynb}.  Data: BTCB-USD hourly OHLCV, March~2024 -- February~2026; USDT-USD hourly peg data.  Libraries: PyTorch~2.x, SHAP~0.51, NumPy, Pandas, Matplotlib, yfinance.}

% ─────────────────────────────────────────────────────────────────────────────
\newpage
\appendix
\section*{Appendix: Full Python Code}

All analysis implemented in Python~3.11. Libraries: NumPy, Pandas, PyTorch~2.x, SHAP~0.51, Matplotlib, yfinance.
Source notebook: \texttt{rl\_crypto\_part3.ipynb}.

\subsection*{Cell 2: All imports}

\begin{lstlisting}[language=Python]
import numpy as np
import pandas as pd
import matplotlib
matplotlib.use('Agg')
import matplotlib.pyplot as plt
import matplotlib.patches as mpatches
import torch
import torch.nn as nn
import torch.optim as optim
import torch.nn.functional as F
from torch.distributions import Categorical
import shap
import warnings
import random
from collections import deque
import os

warnings.filterwarnings('ignore')

# Reproducibility
SEED = 42
np.random.seed(SEED)
torch.manual_seed(SEED)
random.seed(SEED)

# Paths
DATA_PATH = '/Users/ishan/Downloads/BTCB-USD_DataHr.csv'
OUT_DIR = '/Users/ishan/Downloads'

# Hyperparameters
EPISODE_LEN = 168       # one week
GAMMA = 0.99
TRANSACTION_COST = 0.001
STATE_DIM = 13          # 12 features + portfolio fraction
N_ACTIONS = 3           # sell=0, hold=1, buy=2
REINFORCE_LR = 3e-4
BASELINE_LR = 1e-3
DQN_LR = 1e-4
HIDDEN_DIM = 128
REINFORCE_EPISODES = 1200
DQN_EPISODES = 800

print("All imports loaded successfully.")
print(f"PyTorch version: {torch.__version__}")
print(f"SHAP version: {shap.__version__}")
\end{lstlisting}

\subsection*{Cell 3: Data loading + feature engineering}

\begin{lstlisting}[language=Python]
def load_and_engineer(path):
    df = pd.read_csv(path, skiprows=[1, 2], index_col=0, parse_dates=True)
    df.columns = ['Close', 'High', 'Low', 'Open', 'Volume']
    df = df.sort_index().dropna(subset=['Close'])

    close = df['Close']
    high  = df['High']
    low   = df['Low']
    vol   = df['Volume'].replace(0, np.nan).fillna(method='ffill').fillna(1)

    # Returns
    df['ret_1h']   = close.pct_change(1)
    df['ret_4h']   = close.pct_change(4)
    df['ret_24h']  = close.pct_change(24)
    df['ret_168h'] = close.pct_change(168)

    # Volatility (rolling std of 1h returns, annualized proxy)
    df['vol_24h'] = df['ret_1h'].rolling(24).std()

    # RSI 14
    delta = close.diff()
    gain = delta.clip(lower=0).rolling(14).mean()
    loss = (-delta.clip(upper=0)).rolling(14).mean()
    rs   = gain / (loss + 1e-9)
    df['rsi_14'] = 100 - (100 / (1 + rs))

    # Price vs moving averages (normalised)
    ma20 = close.rolling(20).mean()
    ma50 = close.rolling(50).mean()
    df['price_vs_ma20'] = (close - ma20) / (ma20 + 1e-9)
    df['price_vs_ma50'] = (close - ma50) / (ma50 + 1e-9)

    # Time cyclical encoding
    if df.index.tz is not None:
        hour = df.index.tz_convert('UTC').hour
    else:
        hour = df.index.hour
    df['hour_sin'] = np.sin(2 * np.pi * hour / 24)
    df['hour_cos'] = np.cos(2 * np.pi * hour / 24)

    # Log volume
    df['log_vol'] = np.log1p(vol)

    # High-Low range (normalised)
    df['hl_range'] = (high - low) / (close + 1e-9)

    FEATURE_COLS = ['ret_1h', 'ret_4h', 'ret_24h', 'ret_168h',
                    'vol_24h', 'rsi_14', 'price_vs_ma20', 'price_vs_ma50',
                    'hour_sin', 'hour_cos', 'log_vol', 'hl_range']

    df = df.dropna()

    # Market regime labels based on ret_24h
    def label_regime(r):
        if r > 0.01:   return 'bull'
        elif r < -0.01: return 'decline'
        else:           return 'sideways'
    df['regime'] = df['ret_24h'].apply(label_regime)

    # 80/20 split
    n = len(df)
    n_train = int(n * 0.8)
    train_df = df.iloc[:n_train].copy()
    test_df  = df.iloc[n_train:].copy()

    print(f"Total rows after feature engineering: {n}")
    print(f"Train: {len(train_df)}, Test: {len(test_df)}")
    print(f"Features: {FEATURE_COLS}")

    return train_df, test_df, FEATURE_COLS

train_df, test_df, FEATURE_COLS = load_and_engineer(DATA_PATH)

# Normalise features using training statistics
feat_mean = train_df[FEATURE_COLS].mean()
feat_std  = train_df[FEATURE_COLS].std().replace(0, 1)

def normalise(df):
    return (df[FEATURE_COLS] - feat_mean) / feat_std

train_features = normalise(train_df).values.astype(np.float32)
test_features  = normalise(test_df).values.astype(np.float32)
train_prices   = train_df['Close'].values.astype(np.float32)
test_prices    = test_df['Close'].values.astype(np.float32)
train_regime   = train_df['regime'].values
test_regime    = test_df['regime'].values

# Raw feature arrays for SHAP (before normalisation)
train_raw = train_df[FEATURE_COLS].values.astype(np.float32)
test_raw  = test_df[FEATURE_COLS].values.astype(np.float32)

print(f"\nNormalised train feature shape: {train_features.shape}")
print(f"Normalised test  feature shape: {test_features.shape}")
\end{lstlisting}

\subsection*{Cell 4: CryptoEnv base class + TraderEnv subclass with 3 reward functions}

\begin{lstlisting}[language=Python]
class CryptoEnv:
    """
    Base environment for cryptocurrency trading.
    Actions: 0=sell, 1=hold, 2=buy
    State: 12 normalised features + portfolio crypto fraction (0..1)
    """

    def __init__(self, features, prices, episode_len=EPISODE_LEN, tc=TRANSACTION_COST, regime=None):
        self.features    = features          # (T, 12) float32
        self.prices      = prices            # (T,) float32
        self.episode_len = episode_len
        self.tc          = tc
        self.regime      = regime if regime is not None else np.array(['sideways'] * len(prices))
        self.T           = len(prices)
        self.reset()

    # ------------------------------------------------------------------ #
    def reset(self, start=None):
        max_start = self.T - self.episode_len - 1
        if max_start <= 0:
            start = 0
        elif start is None:
            start = np.random.randint(0, max_start)
        self.idx          = start
        self.start_idx    = start
        self.end_idx      = start + self.episode_len
        self.cash         = 1.0
        self.crypto_held  = 0.0
        self.entry_price  = self.prices[self.idx]
        return self._obs()

    def _portfolio_value(self):
        return self.cash + self.crypto_held * self.prices[self.idx]

    def _crypto_fraction(self):
        pv = self._portfolio_value()
        return (self.crypto_held * self.prices[self.idx]) / (pv + 1e-9)

    def _obs(self):
        feat = self.features[self.idx]                    # (12,)
        frac = np.array([self._crypto_fraction()], dtype=np.float32)
        return np.concatenate([feat, frac])               # (13,)

    def _base_reward(self, action):
        """Compute portfolio return after executing action."""
        price_now  = self.prices[self.idx]
        price_next = self.prices[self.idx + 1]
        pv_before  = self._portfolio_value()

        if action == 2:   # buy: invest all cash
            cost = self.cash * (1 - self.tc)
            self.crypto_held += cost / price_now
            self.cash = 0.0
        elif action == 0: # sell: liquidate all crypto
            proceeds = self.crypto_held * price_now * (1 - self.tc)
            self.cash += proceeds
            self.crypto_held = 0.0
        # hold: do nothing

        # Advance time
        self.idx += 1
        pv_after = self._portfolio_value()
        return float(pv_after / (pv_before + 1e-9) - 1.0)

    def step(self, action):
        reward     = self._base_reward(action)
        done       = (self.idx >= self.end_idx)
        obs        = self._obs() if not done else np.zeros(STATE_DIM, dtype=np.float32)
        return obs, reward, done, {}


class TraderEnv(CryptoEnv):
    """
    Extends CryptoEnv with trader-type-specific reward shaping.
    trader_type: 'arbitrageur' | 'manipulator' | 'herder'
    """

    def __init__(self, trader_type, *args, **kwargs):
        super().__init__(*args, **kwargs)
        assert trader_type in ('arbitrageur', 'manipulator', 'herder')
        self.trader_type = trader_type

    def step(self, action):
        # Cache pre-step features for reward shaping
        feat      = self.features[self.idx]    # normalised features
        vol_24h   = float(feat[4])             # index 4
        log_vol   = float(feat[10])            # index 10
        ret_4h    = float(feat[1])             # index 1
        ret_24h   = float(feat[2])             # index 2

        base_reward = self._base_reward(action)

        # Action direction: buy=+1, hold=0, sell=-1
        action_direction = {2: 1.0, 1: 0.0, 0: -1.0}[action]

        if self.trader_type == 'arbitrageur':
            # Risk-adjusted: penalise volatility exposure when holding crypto
            crypto_frac = self._crypto_fraction()
            reward = base_reward - 0.4 * abs(vol_24h) * crypto_frac

        elif self.trader_type == 'manipulator':
            # Front-running: reward trading in direction of large-volume moves
            reward = base_reward + 0.15 * log_vol * action_direction * np.sign(ret_4h)

        elif self.trader_type == 'herder':
            # Momentum: reward aligning action with 24h trend
            reward = base_reward + 0.25 * np.sign(ret_24h) * action_direction

        done = (self.idx >= self.end_idx)
        obs  = self._obs() if not done else np.zeros(STATE_DIM, dtype=np.float32)
        return obs, float(reward), done, {}


# Quick sanity check
env_test = TraderEnv('arbitrageur', train_features, train_prices, regime=train_regime)
s = env_test.reset()
s2, r, done, _ = env_test.step(2)
print(f"Sanity check -- obs shape: {s.shape}, reward: {r:.6f}, done: {done}")
print(f"State dim: {STATE_DIM} = 12 features + 1 portfolio fraction")
\end{lstlisting}

\subsection*{Cell 5: PolicyNet + BaselineNet + REINFORCEAgent}

\begin{lstlisting}[language=Python]
class PolicyNet(nn.Module):
    """
    Two-layer MLP policy network.
    Input: STATE_DIM -> Hidden -> Hidden -> N_ACTIONS (softmax)
    """
    def __init__(self, state_dim=STATE_DIM, hidden=HIDDEN_DIM, n_actions=N_ACTIONS):
        super().__init__()
        self.net = nn.Sequential(
            nn.Linear(state_dim, hidden),
            nn.ReLU(),
            nn.Linear(hidden, hidden),
            nn.ReLU(),
            nn.Linear(hidden, n_actions)
        )

    def forward(self, x):
        return F.softmax(self.net(x), dim=-1)

    def select_action(self, state_np):
        """Sample action and return (action, log_prob)."""
        state = torch.FloatTensor(state_np).unsqueeze(0)
        probs = self.forward(state)
        dist  = Categorical(probs)
        a     = dist.sample()
        return a.item(), dist.log_prob(a)

    def greedy_action(self, state_np):
        """Deterministic greedy action for evaluation."""
        state = torch.FloatTensor(state_np).unsqueeze(0)
        with torch.no_grad():
            probs = self.forward(state)
        return probs.argmax(dim=-1).item()

    def action_probs(self, state_np):
        """Return probability vector as numpy array."""
        state = torch.FloatTensor(state_np).unsqueeze(0)
        with torch.no_grad():
            return self.forward(state).squeeze(0).numpy()


class BaselineNet(nn.Module):
    """
    Two-layer MLP baseline (value) network.
    Input: STATE_DIM -> Hidden -> Hidden -> 1 (state value estimate)
    """
    def __init__(self, state_dim=STATE_DIM, hidden=HIDDEN_DIM):
        super().__init__()
        self.net = nn.Sequential(
            nn.Linear(state_dim, hidden),
            nn.ReLU(),
            nn.Linear(hidden, hidden),
            nn.ReLU(),
            nn.Linear(hidden, 1)
        )

    def forward(self, x):
        return self.net(x).squeeze(-1)


class REINFORCEAgent:
    """
    REINFORCE with learned baseline (variance-reduction).

    Algorithm:
      1. Collect a full episode using the stochastic policy.
      2. Compute discounted returns G_t for each step.
      3. Compute advantages: A_t = G_t - V(s_t)  (baseline subtracted).
      4. Policy loss:  -sum_t( log pi(a_t|s_t) * A_t )
      5. Baseline loss: MSE( V(s_t), G_t )
      6. Update both networks.
    """

    def __init__(self, state_dim=STATE_DIM, hidden=HIDDEN_DIM, n_actions=N_ACTIONS,
                 lr_policy=REINFORCE_LR, lr_baseline=BASELINE_LR, gamma=GAMMA):
        self.policy   = PolicyNet(state_dim, hidden, n_actions)
        self.baseline = BaselineNet(state_dim, hidden)
        self.opt_p    = optim.Adam(self.policy.parameters(),   lr=lr_policy)
        self.opt_b    = optim.Adam(self.baseline.parameters(), lr=lr_baseline)
        self.gamma    = gamma

    def collect_episode(self, env):
        """Run one episode, return (states, actions, rewards, log_probs)."""
        states, actions, rewards, log_probs = [], [], [], []
        state = env.reset()
        done  = False
        while not done:
            a, lp    = self.policy.select_action(state)
            ns, r, done, _ = env.step(a)
            states.append(state)
            actions.append(a)
            rewards.append(r)
            log_probs.append(lp)
            state = ns
        return states, actions, rewards, log_probs

    def _discounted_returns(self, rewards):
        G, running = [], 0.0
        for r in reversed(rewards):
            running = r + self.gamma * running
            G.insert(0, running)
        return torch.FloatTensor(G)

    def update(self, states, rewards, log_probs):
        G = self._discounted_returns(rewards)

        states_t = torch.FloatTensor(np.array(states))
        V        = self.baseline(states_t)

        # Normalise advantages for training stability
        A = G - V.detach()
        A = (A - A.mean()) / (A.std() + 1e-8)

        # Policy (actor) loss
        log_probs_t = torch.stack(log_probs)
        policy_loss = -(log_probs_t * A).sum()

        self.opt_p.zero_grad()
        policy_loss.backward()
        nn.utils.clip_grad_norm_(self.policy.parameters(), 1.0)
        self.opt_p.step()

        # Baseline (critic) MSE loss
        baseline_loss = F.mse_loss(V, G)
        self.opt_b.zero_grad()
        baseline_loss.backward()
        self.opt_b.step()

        return float(sum(rewards)), float(policy_loss), float(baseline_loss)


# Test instantiation
agent_test = REINFORCEAgent()
print(f"PolicyNet parameters:   {sum(p.numel() for p in agent_test.policy.parameters()):,}")
print(f"BaselineNet parameters: {sum(p.numel() for p in agent_test.baseline.parameters()):,}")
print("REINFORCEAgent created successfully.")
\end{lstlisting}

\subsection*{Cell 6: Train REINFORCE for all 3 trader types}

\begin{lstlisting}[language=Python]
import time

def train_reinforce(trader_type, n_episodes=REINFORCE_EPISODES, print_every=200):
    """Train a REINFORCE agent for a given trader type."""
    env   = TraderEnv(trader_type, train_features, train_prices, regime=train_regime)
    agent = REINFORCEAgent()
    ep_rewards = []
    t0 = time.time()

    for ep in range(1, n_episodes + 1):
        states, actions, rewards, log_probs = agent.collect_episode(env)
        ep_r, p_loss, b_loss = agent.update(states, rewards, log_probs)
        ep_rewards.append(ep_r)

        if ep % print_every == 0:
            avg = np.mean(ep_rewards[-print_every:])
            elapsed = time.time() - t0
            print(f"  [{trader_type:12s}] Ep {ep:4d}/{n_episodes}  "
                  f"AvgReward(last{print_every}): {avg:+.5f}  "
                  f"PolicyLoss: {p_loss:+.4f}  "
                  f"Elapsed: {elapsed:.1f}s")

    total_time = time.time() - t0
    print(f"  [{trader_type:12s}] Training complete. Total: {total_time:.1f}s")
    return agent, ep_rewards


print("=" * 65)
print("Training REINFORCE agents for 3 trader types...")
print("=" * 65)

agents = {}
reward_histories = {}

for ttype in ['arbitrageur', 'manipulator', 'herder']:
    print(f"\n--- {ttype.upper()} ---")
    agent, history = train_reinforce(ttype, n_episodes=REINFORCE_EPISODES)
    agents[ttype]         = agent
    reward_histories[ttype] = history

print("\nAll REINFORCE agents trained.")
\end{lstlisting}

\subsection*{Cell 7: Backtest all 3 REINFORCE agents + Buy \& Hold on test set}

\begin{lstlisting}[language=Python]
def backtest(policy_net, features, prices, regime=None, tc=TRANSACTION_COST):
    """
    Run deterministic greedy backtest on the full provided test data.
    Returns dict with metrics and portfolio value series.
    """
    if regime is None:
        regime = np.array(['sideways'] * len(prices))

    T            = len(prices)
    cash         = 1.0
    crypto_held  = 0.0
    pv_series    = [1.0]
    returns_list = []

    for t in range(T - 1):
        pv_before = cash + crypto_held * prices[t]
        frac      = (crypto_held * prices[t]) / (pv_before + 1e-9)
        feat_vec  = features[t]
        state     = np.concatenate([feat_vec, [frac]]).astype(np.float32)
        action    = policy_net.greedy_action(state)

        if action == 2:   # buy
            cost         = cash * (1 - tc)
            crypto_held += cost / prices[t]
            cash         = 0.0
        elif action == 0: # sell
            proceeds     = crypto_held * prices[t] * (1 - tc)
            cash        += proceeds
            crypto_held  = 0.0

        pv_after = cash + crypto_held * prices[t + 1]
        ret      = pv_after / (pv_before + 1e-9) - 1.0
        returns_list.append(ret)
        pv_series.append(pv_after)

    pv_arr     = np.array(pv_series)
    ret_arr    = np.array(returns_list)
    total_ret  = pv_arr[-1] - 1.0
    sharpe     = (ret_arr.mean() / (ret_arr.std() + 1e-9)) * np.sqrt(8760)
    rolling_max = np.maximum.accumulate(pv_arr)
    drawdowns   = (pv_arr - rolling_max) / (rolling_max + 1e-9)
    max_dd      = drawdowns.min()

    return {
        'Total Return': total_ret,
        'Sharpe Ratio': sharpe,
        'Max Drawdown': max_dd,
        'Final Portfolio': pv_arr[-1],
        'pv_series': pv_arr
    }

def buy_and_hold(prices):
    """Buy at start, hold until end."""
    pv_series = prices / prices[0]
    returns   = np.diff(pv_series) / (pv_series[:-1] + 1e-9)
    total_ret = pv_series[-1] - 1.0
    sharpe    = (returns.mean() / (returns.std() + 1e-9)) * np.sqrt(8760)
    rolling_max = np.maximum.accumulate(pv_series)
    drawdowns   = (pv_series - rolling_max) / (rolling_max + 1e-9)
    max_dd      = drawdowns.min()
    return {
        'Total Return': total_ret,
        'Sharpe Ratio': sharpe,
        'Max Drawdown': max_dd,
        'Final Portfolio': pv_series[-1],
        'pv_series': pv_series
    }

print("Running backtests on test set...")
backtest_results = {}

for ttype in ['arbitrageur', 'manipulator', 'herder']:
    res = backtest(agents[ttype].policy, test_features, test_prices, regime=test_regime)
    backtest_results[f'REINFORCE-{ttype[:4]}'] = res
    print(f"  REINFORCE-{ttype[:4]:4s}  Return: {res['Total Return']:+.4f}  "
          f"Sharpe: {res['Sharpe Ratio']:+.3f}  MaxDD: {res['Max Drawdown']:.4f}  "
          f"FinalPV: {res['Final Portfolio']:.4f}")

bh = buy_and_hold(test_prices)
backtest_results['Buy&Hold'] = bh
print(f"  Buy&Hold      Return: {bh['Total Return']:+.4f}  "
      f"Sharpe: {bh['Sharpe Ratio']:+.3f}  MaxDD: {bh['Max Drawdown']:.4f}  "
      f"FinalPV: {bh['Final Portfolio']:.4f}")

print("\nBacktest complete.")
\end{lstlisting}

\subsection*{Cell 8: DQN Comparison --- Double DQN with 3-layer MLP}

\begin{lstlisting}[language=Python]
class QNetwork(nn.Module):
    """3-layer MLP Q-network (same architecture as Part 2)."""
    def __init__(self, state_dim=STATE_DIM, hidden=HIDDEN_DIM, n_actions=N_ACTIONS):
        super().__init__()
        self.net = nn.Sequential(
            nn.Linear(state_dim, hidden),
            nn.ReLU(),
            nn.Linear(hidden, hidden),
            nn.ReLU(),
            nn.Linear(hidden, hidden // 2),
            nn.ReLU(),
            nn.Linear(hidden // 2, n_actions)
        )

    def forward(self, x):
        return self.net(x)


class ReplayBuffer:
    def __init__(self, capacity=50_000):
        self.buf = deque(maxlen=capacity)

    def push(self, s, a, r, ns, done):
        self.buf.append((s, a, r, ns, done))

    def sample(self, batch_size):
        batch = random.sample(self.buf, batch_size)
        s, a, r, ns, d = zip(*batch)
        return (torch.FloatTensor(np.array(s)),
                torch.LongTensor(a),
                torch.FloatTensor(r),
                torch.FloatTensor(np.array(ns)),
                torch.FloatTensor(d))

    def __len__(self):
        return len(self.buf)


class DQNAgent:
    """Double DQN agent replicating Part 2 architecture."""

    def __init__(self, state_dim=STATE_DIM, hidden=HIDDEN_DIM, n_actions=N_ACTIONS,
                 lr=DQN_LR, gamma=GAMMA, epsilon_start=1.0, epsilon_end=0.05,
                 epsilon_decay=0.995, batch_size=64, target_update=10):
        self.q_net      = QNetwork(state_dim, hidden, n_actions)
        self.target_net = QNetwork(state_dim, hidden, n_actions)
        self.target_net.load_state_dict(self.q_net.state_dict())
        self.target_net.eval()
        self.opt         = optim.Adam(self.q_net.parameters(), lr=lr)
        self.buffer      = ReplayBuffer()
        self.gamma       = gamma
        self.epsilon     = epsilon_start
        self.eps_end     = epsilon_end
        self.eps_decay   = epsilon_decay
        self.batch_size  = batch_size
        self.target_upd  = target_update
        self.n_actions   = n_actions
        self.step_count  = 0

    def select_action(self, state_np):
        if random.random() < self.epsilon:
            return random.randint(0, self.n_actions - 1)
        with torch.no_grad():
            q = self.q_net(torch.FloatTensor(state_np).unsqueeze(0))
        return q.argmax().item()

    def greedy_action(self, state_np):
        with torch.no_grad():
            q = self.q_net(torch.FloatTensor(state_np).unsqueeze(0))
        return q.argmax().item()

    def push(self, s, a, r, ns, done):
        self.buffer.push(s, a, r, ns, float(done))

    def learn(self):
        if len(self.buffer) < self.batch_size:
            return
        s, a, r, ns, d = self.buffer.sample(self.batch_size)

        # Double DQN
        with torch.no_grad():
            best_a   = self.q_net(ns).argmax(1, keepdim=True)
            q_target = r + self.gamma * (1 - d) * self.target_net(ns).gather(1, best_a).squeeze(1)

        q_pred = self.q_net(s).gather(1, a.unsqueeze(1)).squeeze(1)
        loss   = F.mse_loss(q_pred, q_target)
        self.opt.zero_grad()
        loss.backward()
        nn.utils.clip_grad_norm_(self.q_net.parameters(), 1.0)
        self.opt.step()

        self.step_count += 1
        if self.step_count % self.target_upd == 0:
            self.target_net.load_state_dict(self.q_net.state_dict())

        self.epsilon = max(self.eps_end, self.epsilon * self.eps_decay)


def train_dqn(n_episodes=DQN_EPISODES, print_every=200):
    """Train DQN on the general (base) CryptoEnv."""
    env    = CryptoEnv(train_features, train_prices, regime=train_regime)
    agent  = DQNAgent()
    ep_rewards = []
    t0 = time.time()

    for ep in range(1, n_episodes + 1):
        state  = env.reset()
        done   = False
        ep_r   = 0.0
        while not done:
            a              = agent.select_action(state)
            ns, r, done, _ = env.step(a)
            agent.push(state, a, r, ns, done)
            agent.learn()
            ep_r  += r
            state  = ns
        ep_rewards.append(ep_r)

        if ep % print_every == 0:
            avg     = np.mean(ep_rewards[-print_every:])
            elapsed = time.time() - t0
            print(f"  [DQN] Ep {ep:4d}/{n_episodes}  "
                  f"AvgReward(last{print_every}): {avg:+.5f}  "
                  f"Epsilon: {agent.epsilon:.3f}  Elapsed: {elapsed:.1f}s")

    print(f"  [DQN] Training complete. Total: {time.time()-t0:.1f}s")
    return agent, ep_rewards


print("Training DQN agent (Double DQN, 3-layer MLP, 800 episodes)...")
dqn_agent, dqn_history = train_dqn(n_episodes=DQN_EPISODES)

# Backtest DQN
class DQNPolicyWrapper:
    """Wrap DQN to match PolicyNet interface expected by backtest()."""
    def __init__(self, dqn):
        self.dqn = dqn
    def greedy_action(self, state_np):
        return self.dqn.greedy_action(state_np)

dqn_res = backtest(DQNPolicyWrapper(dqn_agent), test_features, test_prices, regime=test_regime)
backtest_results['DQN'] = dqn_res

print(f"\n  DQN           Return: {dqn_res['Total Return']:+.4f}  "
      f"Sharpe: {dqn_res['Sharpe Ratio']:+.3f}  "
      f"MaxDD: {dqn_res['Max Drawdown']:.4f}  "
      f"FinalPV: {dqn_res['Final Portfolio']:.4f}")

print("\n" + "=" * 75)
print("COMPARISON TABLE: REINFORCE (3 trader types) vs DQN vs Buy&Hold")
print("=" * 75)
print(f"{'Agent':<22} {'Total Return':>14} {'Sharpe':>8} {'Max Drawdown':>14} {'Final PV':>10}")
print("-" * 75)
for name, res in backtest_results.items():
    print(f"{name:<22} {res['Total Return']:>14.4f} {res['Sharpe Ratio']:>8.3f} "
          f"{res['Max Drawdown']:>14.4f} {res['Final Portfolio']:>10.4f}")
print("=" * 75)
\end{lstlisting}

\subsection*{Cell 9: SHAP Analysis --- Kernel SHAP for all 3 trader types}

\begin{lstlisting}[language=Python]
def collect_states_greedy(policy_net, features, prices, n_steps=200):
    """
    Run greedy policy on test env to collect states.
    Returns array of shape (n_steps, 12) using the raw (unnormalised) features
    plus portfolio fraction.
    We return two arrays:
      - states_norm:  (n_steps, STATE_DIM) normalised features + frac
      - states_raw12: (n_steps, 12) raw (unnormalised) feature values  (for SHAP display)
    """
    T = len(prices)
    cash         = 1.0
    crypto_held  = 0.0
    states_norm  = []
    states_raw12 = []

    for t in range(min(n_steps, T - 1)):
        pv   = cash + crypto_held * prices[t]
        frac = (crypto_held * prices[t]) / (pv + 1e-9)
        feat_norm = features[t]
        state     = np.concatenate([feat_norm, [frac]]).astype(np.float32)
        states_norm.append(state)
        states_raw12.append(feat_norm)          # normalised features -- used as proxy

        action = policy_net.greedy_action(state)
        if action == 2:
            cost         = cash * (1 - TRANSACTION_COST)
            crypto_held += cost / prices[t]
            cash         = 0.0
        elif action == 0:
            proceeds     = crypto_held * prices[t] * (1 - TRANSACTION_COST)
            cash        += proceeds
            crypto_held  = 0.0

    return np.array(states_norm), np.array(states_raw12)


def shap_predict_buy_prob(model, X):
    """
    SHAP-compatible prediction function.
    Takes (n, STATE_DIM) numpy array, returns (n,) Buy-probability.
    We pad the 12 or 13-dim input appropriately.
    """
    n = X.shape[0]
    # X has 12 features (background only has feat columns, not frac)
    # We append a zero portfolio fraction to make STATE_DIM=13
    if X.shape[1] == 12:
        frac = np.zeros((n, 1), dtype=np.float32)
        X_full = np.concatenate([X, frac], axis=1).astype(np.float32)
    else:
        X_full = X.astype(np.float32)

    with torch.no_grad():
        t = torch.FloatTensor(X_full)
        probs = model(t).numpy()   # (n, 3)
    return probs[:, 2]             # Buy probability (action=2)


print("Running SHAP KernelExplainer for each trader type...")
print("(Background: 30 states, Test: 100 states, nsamples=80)")
print()

shap_results = {}   # {trader_type: {'shap_vals': array, 'mean_abs': array}}
N_BACKGROUND = 30
N_TEST       = 100
N_SHAP_SAMP  = 80

for ttype in ['arbitrageur', 'manipulator', 'herder']:
    print(f"--- {ttype.upper()} ---")
    policy = agents[ttype].policy

    # Collect states
    states_norm, states_raw12 = collect_states_greedy(policy, test_features, test_prices, n_steps=200)
    print(f"  Collected {len(states_norm)} states from greedy rollout")

    bg_raw  = states_raw12[:N_BACKGROUND]       # (30, 12)
    tst_raw = states_raw12[N_BACKGROUND:N_BACKGROUND + N_TEST]  # (100, 12)

    def pred_fn(X):
        return shap_predict_buy_prob(policy, X)

    explainer = shap.KernelExplainer(pred_fn, bg_raw)
    shap_vals = explainer.shap_values(tst_raw, nsamples=N_SHAP_SAMP)
    # shap_vals shape: (100, 12)

    mean_abs = np.abs(shap_vals).mean(axis=0)   # (12,)

    # Ranked feature importance
    ranked_idx = np.argsort(mean_abs)[::-1]
    print(f"  Feature importance (ranked by mean |SHAP|):")
    for rank, idx in enumerate(ranked_idx):
        print(f"    #{rank+1:2d}  {FEATURE_COLS[idx]:>18s}  mean|SHAP|= {mean_abs[idx]:.6f}")

    shap_results[ttype] = {
        'shap_vals': shap_vals,
        'mean_abs':  mean_abs,
        'ranked_idx': ranked_idx
    }
    print()

print("SHAP analysis complete.")

# Summary: Top 3 features per trader type
print("\n" + "=" * 55)
print("TOP 3 FEATURES PER TRADER TYPE (mean |SHAP| for Buy prob)")
print("=" * 55)
for ttype in ['arbitrageur', 'manipulator', 'herder']:
    res = shap_results[ttype]
    print(f"\n{ttype.upper()}:")
    for rank in range(3):
        idx = res['ranked_idx'][rank]
        print(f"  #{rank+1}: {FEATURE_COLS[idx]:>18s}  {res['mean_abs'][idx]:.6f}")
print("=" * 55)
\end{lstlisting}

\subsection*{Cell 10: Market Simulation --- Crisis Frequency, Stablecoin Stability, Market Efficiency}

\begin{lstlisting}[language=Python]
import yfinance as yf

# Download USDT-USD for stablecoin stability (actual peg deviation data)
_start = ftest.index[0].strftime('%Y-%m-%d')
_end   = ftest.index[-1].strftime('%Y-%m-%d')
usdt_raw = yf.download('USDT-USD', start=_start, end=_end, interval='1h',
                        auto_adjust=True, progress=False)
usdt_close = usdt_raw['Close']
if isinstance(usdt_close, pd.DataFrame): usdt_close = usdt_close.iloc[:,0]
if usdt_close.index.tzinfo: usdt_close.index = usdt_close.index.tz_convert(None)
usdt_peg_dev = (usdt_close - 1.0).abs() * 100   # % deviation from $1.00 peg
usdt_arr = usdt_peg_dev.reindex(ftest.index, method='nearest',
                                 tolerance=pd.Timedelta('2h')).fillna(usdt_peg_dev.mean()).values

print(f"USDT aligned: {(~np.isnan(usdt_arr)).sum()} hours, mean peg dev: {usdt_arr.mean():.4f}%")

feat_test   = ftest.values.astype(np.float32)
price_test  = ptest.values.astype(np.float32)
regime_arr  = rtest.values
SIM_STEPS   = min(500, len(feat_test) - 2)

def a2sig(a): return {1: 1., 0: 0., 2: -1.}[a]

def simulate_market(w_arb, w_man, w_hrd):
    """
    Simulate 500 hours on test data under a given trader-type composition.
    Returns: crisis_frequency, stablecoin_stability (USDT-based), market_efficiency.
    """
    agents  = [(arb_agent, w_arb), (man_agent, w_man), (hrd_agent, w_hrd)]
    crypto_frac = 0.0
    crisis_steps = 0
    sim_rets, signals, stab_vals = [], [], []
    price = price_test[0]

    for t in range(SIM_STEPS):
        state = np.append(feat_test[t], crypto_frac).astype(np.float32)
        agg_sig = sum(w * a2sig(ag.act(state, greedy=True)) for ag, w in agents)
        signals.append(agg_sig)

        # Price with micro market-impact (0.05% per unit signal)
        impact   = agg_sig * 0.0005
        new_price = price * (1 + feat_test[t, FEAT_COLS.index('ret_1h')] + impact)
        ret = (new_price - price) / (price + 1e-9)
        sim_rets.append(ret)
        price = new_price

        # Update portfolio exposure based on aggregate action
        if   agg_sig >  0.3: crypto_frac = min(crypto_frac + 0.1, 1.0)
        elif agg_sig < -0.3: crypto_frac = max(crypto_frac - 0.1, 0.0)

        # Crisis: confirmed decline regime + strong collective sell signal
        if regime_arr[t] == 'decline' and agg_sig < -0.3:
            crisis_steps += 1

        stab_vals.append(usdt_arr[t])

    sigs = np.array(signals)
    # Stablecoin stability: USDT peg deviation DURING high-stress periods
    stressed = np.abs(sigs) > 0.3
    usdt_stress = np.array(stab_vals)[stressed].mean() if stressed.sum() > 5 else np.array(stab_vals).mean()
    stablecoin_stability = 1.0 / (usdt_stress + 1e-6)   # higher = more stable

    sim_rets = np.array(sim_rets)
    lag1_ac  = np.corrcoef(sim_rets[:-1], sim_rets[1:])[0, 1] if len(sim_rets) > 2 else 0.
    market_efficiency = 1.0 - abs(lag1_ac)

    return {
        'crisis_freq':           crisis_steps / SIM_STEPS,
        'stablecoin_stability':  stablecoin_stability,
        'market_efficiency':     market_efficiency
    }

compositions = [
    (1.0, 0.0, 0.0), (0.0, 1.0, 0.0), (0.0, 0.0, 1.0),
    (0.33, 0.33, 0.34), (0.6, 0.2, 0.2), (0.2, 0.6, 0.2), (0.2, 0.2, 0.6)
]
comp_labels = ["All Arb","All Manip","All Herd","Equal","Arb-Dom","Manip-Dom","Herd-Dom"]

sim_results = {}
for (wa, wm, wh), lb in zip(compositions, comp_labels):
    sim_results[lb] = simulate_market(wa, wm, wh)

print("\nMarket Simulation Results:")
print(f"{'Composition':12s}  {'Crisis Freq':>12s}  {'USDT Stability':>15s}  {'Mkt Efficiency':>14s}")
for lb in comp_labels:
    r = sim_results[lb]
    print(f"{lb:12s}  {r['crisis_freq']:>12.4f}  {r['stablecoin_stability']:>15.4f}  {r['market_efficiency']:>14.4f}")
\end{lstlisting}

\subsection*{Cell 11: Three-panel market metrics figure}

\begin{lstlisting}[language=Python]
bar_colors = {
    "All Arb":    '#2166AC',   # blue
    "All Manip":  '#D6604D',   # red
    "All Herd":   '#F4A582',   # orange
    "Equal":      '#762A83',   # purple
    "Arb-Dom":    '#4393C3',   # light blue
    "Manip-Dom":  '#B2182B',   # dark red
    "Herd-Dom":   '#E08214',   # dark orange
}

crisis_vals  = [sim_results[l]['crisis_freq']      for l in comp_labels]
stab_vals    = [sim_results[l]['price_stability']   for l in comp_labels]
eff_vals     = [sim_results[l]['market_efficiency'] for l in comp_labels]
colors       = [bar_colors[l] for l in comp_labels]
x            = np.arange(len(comp_labels))

equal_idx = comp_labels.index("Equal")

fig, axes = plt.subplots(1, 3, figsize=(18, 6))
fig.suptitle("Market Composition vs. Stability Metrics\n(REINFORCE Multi-Trader Simulation)",
             fontsize=14, fontweight='bold', y=1.02)

# Panel 1: Crisis Frequency
ax = axes[0]
bars = ax.bar(x, crisis_vals, color=colors, edgecolor='black', linewidth=0.7, alpha=0.85)
ax.axhline(y=crisis_vals[equal_idx], color='purple', linestyle='--', linewidth=1.8,
           label=f'Equal mix baseline ({crisis_vals[equal_idx]:.3f})')
ax.set_xticks(x); ax.set_xticklabels(comp_labels, rotation=30, ha='right', fontsize=9)
ax.set_ylabel("Crisis Frequency (fraction of steps)")
ax.set_title("Crisis Frequency\n(lower is better)", fontsize=11)
ax.legend(fontsize=8)
ax.set_ylim(0, max(crisis_vals) * 1.35 if max(crisis_vals) > 0 else 0.1)
for bar, v in zip(bars, crisis_vals):
    ax.text(bar.get_x() + bar.get_width()/2., bar.get_height() + 0.001,
            f'{v:.3f}', ha='center', va='bottom', fontsize=8)

# Panel 2: Price Stability
ax = axes[1]
bars = ax.bar(x, stab_vals, color=colors, edgecolor='black', linewidth=0.7, alpha=0.85)
ax.axhline(y=stab_vals[equal_idx], color='purple', linestyle='--', linewidth=1.8,
           label=f'Equal mix baseline ({stab_vals[equal_idx]:.2f})')
ax.set_xticks(x); ax.set_xticklabels(comp_labels, rotation=30, ha='right', fontsize=9)
ax.set_ylabel("Price Stability (1 / std of agg. signals)")
ax.set_title("Price Stability\n(higher is better)", fontsize=11)
ax.legend(fontsize=8)
ax.set_ylim(0, max(stab_vals) * 1.2)
for bar, v in zip(bars, stab_vals):
    ax.text(bar.get_x() + bar.get_width()/2., bar.get_height() + 0.02,
            f'{v:.2f}', ha='center', va='bottom', fontsize=8)

# Panel 3: Market Efficiency
ax = axes[2]
bars = ax.bar(x, eff_vals, color=colors, edgecolor='black', linewidth=0.7, alpha=0.85)
ax.axhline(y=eff_vals[equal_idx], color='purple', linestyle='--', linewidth=1.8,
           label=f'Equal mix baseline ({eff_vals[equal_idx]:.3f})')
ax.set_xticks(x); ax.set_xticklabels(comp_labels, rotation=30, ha='right', fontsize=9)
ax.set_ylabel("Market Efficiency (1 - |lag-1 autocorr|)")
ax.set_title("Market Efficiency\n(higher is better)", fontsize=11)
ax.legend(fontsize=8)
ax.set_ylim(0, 1.15)
for bar, v in zip(bars, eff_vals):
    ax.text(bar.get_x() + bar.get_width()/2., bar.get_height() + 0.005,
            f'{v:.3f}', ha='center', va='bottom', fontsize=8)

# Shared legend for trader types
legend_patches = [
    mpatches.Patch(color='#2166AC', label='Arbitrageur-heavy'),
    mpatches.Patch(color='#D6604D', label='Manipulator-heavy'),
    mpatches.Patch(color='#F4A582', label='Herder-heavy'),
    mpatches.Patch(color='#762A83', label='Mixed'),
]
fig.legend(handles=legend_patches, loc='lower center', ncol=4,
           bbox_to_anchor=(0.5, -0.08), fontsize=9, frameon=True)

plt.tight_layout()
out_path = os.path.join(OUT_DIR, 'part3_market_metrics.png')
plt.savefig(out_path, dpi=150, bbox_inches='tight')
plt.show()
print(f"Saved: {out_path}")
\end{lstlisting}

\subsection*{Cell 12: SHAP Summary Bar Charts (3 subplots, one per trader type)}

\begin{lstlisting}[language=Python]
fig, axes = plt.subplots(1, 3, figsize=(18, 6))
fig.suptitle("SHAP Feature Importance for Buy-Action Probability\n(Mean |SHAP| across 100 test states)",
             fontsize=13, fontweight='bold')

trader_colors = {
    'arbitrageur': '#2166AC',
    'manipulator': '#D6604D',
    'herder':      '#E08214',
}

for ax, ttype in zip(axes, ['arbitrageur', 'manipulator', 'herder']):
    res       = shap_results[ttype]
    mean_abs  = res['mean_abs']   # (12,)
    ranked    = res['ranked_idx']

    # Sort features by mean |SHAP| descending
    sorted_vals  = mean_abs[ranked]
    sorted_names = [FEATURE_COLS[i] for i in ranked]

    ypos = np.arange(len(FEATURE_COLS))
    ax.barh(ypos, sorted_vals, color=trader_colors[ttype], alpha=0.85,
            edgecolor='black', linewidth=0.5)
    ax.set_yticks(ypos)
    ax.set_yticklabels(sorted_names, fontsize=9)
    ax.invert_yaxis()   # top = most important
    ax.set_xlabel("Mean |SHAP| value")
    ax.set_title(f"{ttype.capitalize()}\n(top: {sorted_names[0]})", fontsize=11)

    # Annotate top 3
    for i in range(min(3, len(sorted_vals))):
        ax.text(sorted_vals[i] + sorted_vals[0] * 0.01,
                ypos[i], f'{sorted_vals[i]:.5f}',
                va='center', fontsize=7.5, color='black')

plt.tight_layout()
shap_out = os.path.join(OUT_DIR, 'part3_shap.png')
plt.savefig(shap_out, dpi=150, bbox_inches='tight')
plt.show()
print(f"Saved: {shap_out}")
\end{lstlisting}

\subsection*{Cell 13: Training Curves --- 100-episode rolling average for all 3 REINFORCE types}

\begin{lstlisting}[language=Python]
fig, ax = plt.subplots(figsize=(12, 5))

color_map = {'arbitrageur': '#2166AC', 'manipulator': '#D6604D', 'herder': '#E08214'}
roll_window = 100

for ttype, history in reward_histories.items():
    x_eps  = np.arange(1, len(history) + 1)
    raw    = np.array(history)
    # Rolling average using pandas for clean handling of edges
    rolled = pd.Series(raw).rolling(window=roll_window, min_periods=1).mean().values

    ax.plot(x_eps, raw, alpha=0.15, color=color_map[ttype], linewidth=0.6)
    ax.plot(x_eps, rolled, color=color_map[ttype], linewidth=2.0,
            label=f'{ttype.capitalize()} ({roll_window}-ep MA)')

ax.set_xlabel("Episode", fontsize=11)
ax.set_ylabel("Cumulative Episode Reward", fontsize=11)
ax.set_title("REINFORCE Training Curves by Trader Type\n"
             "(light: raw, solid: 100-episode rolling average)", fontsize=12)
ax.legend(fontsize=10)
ax.grid(True, alpha=0.3)
ax.axhline(y=0, color='black', linewidth=0.8, linestyle='--', alpha=0.5)

plt.tight_layout()
training_out = os.path.join(OUT_DIR, 'part3_training.png')
plt.savefig(training_out, dpi=150, bbox_inches='tight')
plt.show()
print(f"Saved: {training_out}")
\end{lstlisting}

\subsection*{Cell 14: Summary Results Table (formatted DataFrame)}

\begin{lstlisting}[language=Python]
metrics_data = []
for name, res in backtest_results.items():
    metrics_data.append({
        'Agent':          name,
        'Total Return':   f"{res['Total Return']:+.4f}",
        'Sharpe Ratio':   f"{res['Sharpe Ratio']:+.3f}",
        'Max Drawdown':   f"{res['Max Drawdown']:.4f}",
        'Final Portfolio':f"{res['Final Portfolio']:.4f}",
    })

summary_df = pd.DataFrame(metrics_data).set_index('Agent')

print("=" * 68)
print("SUMMARY: All Agents -- Backtest Performance on Test Set")
print("=" * 68)
print(summary_df.to_string())
print("=" * 68)
print()

# Also show raw numeric summary
print("\nDetailed numeric summary:")
metrics_num = []
for name, res in backtest_results.items():
    metrics_num.append({
        'Agent':          name,
        'Total Return':   res['Total Return'],
        'Sharpe Ratio':   res['Sharpe Ratio'],
        'Max Drawdown':   res['Max Drawdown'],
        'Final Portfolio':res['Final Portfolio'],
    })
df_num = pd.DataFrame(metrics_num).set_index('Agent')
print(df_num.round(4).to_string())

# Best agent by Sharpe
best_sharpe = df_num['Sharpe Ratio'].idxmax()
best_return = df_num['Total Return'].idxmax()
print(f"\nBest Sharpe:       {best_sharpe}  ({df_num.loc[best_sharpe,'Sharpe Ratio']:.3f})")
print(f"Best Total Return: {best_return}  ({df_num.loc[best_return,'Total Return']:.4f})")
\end{lstlisting}

\end{document}
